\documentclass[11pt]{article}

\usepackage[margin=1in]{geometry}
\usepackage{amsmath,amssymb}
\usepackage{array}
\usepackage{enumitem}

\title{Lecture 5 Notes: \\ Logicism: Is Mathematics (Just) Logic?}
\author{Philosophy of Mathematics}
\date{}

\begin{document}

\maketitle

\section*{Central Theme}

Logicism is the view that mathematics, or at least a significant part of mathematics, can be
reduced to logic.

\[
\boxed{
\text{Mathematics} = \text{logic} + \text{definitions?}
}
\]

The guiding question:

\[
\boxed{
\text{Can mathematics be analytic rather than synthetic?}
}
\]

Logicism offers a third route after Kant and Mill.

\[
\begin{array}{c|c|c}
\textbf{Kant} & \textbf{Mill} & \textbf{Logicism} \\
\hline
\text{synthetic a priori} & \text{synthetic a posteriori} & \text{analytic a priori} \\
\text{intuition} & \text{experience} & \text{logic and definitions}
\end{array}
\]

\section*{1. The Logicist Alternative}

Kant held that mathematics is synthetic a priori. Mathematical knowledge is not derived
from experience, but it is also not merely conceptual analysis.

Mill took the opposite route. Mathematics is synthetic, but empirical. It is grounded in
experience and induction.

Logicism proposes another answer:

\[
\text{Mathematics is analytic.}
\]

Mathematical truths are knowable a priori because they can be derived from logic and
definitions alone.

\subsection*{Teaching Point}

Logicism tries to preserve the necessity and a priority of mathematics without Kantian
intuition and without Millian empiricism.

\subsection*{Whiteboard}

\[
\text{Kant: intuition}
\]

\[
\text{Mill: experience}
\]

\[
\text{Logicism: logic + definitions}
\]

\section*{2. Frege's New Conception of Analyticity}

For Kant, a proposition is analytic if the predicate concept is contained in the subject
concept.

Frege changes the issue. For Frege, analyticity concerns the ultimate ground of
justification.

\[
\boxed{
\text{A truth is analytic if it is provable from logic and definitions alone.}
}
\]

So the question is not whether \(7+5=12\) is contained in the concepts of \(7\), \(5\), \(+\),
and \(12\). The question is whether arithmetic can be derived from logical laws and
definitions.

\subsection*{Teaching Point}

Frege turns analyticity from a theory of concept-containment into a theory of proof.

\subsection*{Whiteboard}

\[
\text{Kant: analytic} = \text{concept containment}
\]

\[
\text{Frege: analytic} = \text{proof from logic and definitions}
\]

\section*{3. Frege's Program}

Frege's project is:

\[
\boxed{
\text{Define number logically, then derive arithmetic logically.}
}
\]

This makes arithmetic a test case for logicism.

Frege is a realist about mathematics. He thinks numbers are genuine objects and that
mathematical statements have objective truth-values.

\[
\begin{array}{c|c}
\textbf{Realism in ontology} & \text{Numbers exist as objects.} \\
\textbf{Realism in truth-value} & \text{Arithmetic statements are objectively true or false.}
\end{array}
\]

\subsection*{Teaching Point}

Frege is not trying to eliminate numbers. He is trying to show that numbers are logical
objects.

\subsection*{Whiteboard}

\[
\text{Number} \rightarrow \text{logical object}
\]

\[
\text{Arithmetic} \rightarrow \text{logical science}
\]

\section*{4. Equinumerosity}

Frege begins with a basic fact about counting.

Two concepts are equinumerous if the objects falling under one can be put into one-to-one
correspondence with the objects falling under the other.

Examples:

\[
\text{napkins} \sim \text{plates}
\]

if there is exactly one napkin for each plate.

\[
\text{husbands} \sim \text{wives}
\]

in a monogamous society, by definition.

\subsection*{Teaching Point}

Counting can be understood through one-to-one correspondence before we explicitly
introduce numbers.

\subsection*{Whiteboard}

\[
F \sim G
\quad \text{iff}
\quad
F\text{s and }G\text{s can be paired one-to-one.}
\]

\[
\text{same number}
=
\text{equinumerous}
\]

\section*{5. Hume's Principle}

Frege's central principle is now called Hume's Principle.

\[
\boxed{
\#F = \#G
\quad \text{iff}
\quad
F \sim G
}
\]

In words:

\[
\text{The number of }F = \text{the number of }G
\]

if and only if

\[
F \text{ and } G \text{ are equinumerous.}
\]

This principle connects number-identity with one-to-one correspondence.

\subsection*{Teaching Point}

Hume's Principle is the heart of Fregean logicism. It gives identity conditions for numbers
in terms of equinumerosity.

\subsection*{Whiteboard}

\[
\#F = \#G
\Longleftrightarrow
F \sim G
\]

\[
\text{number identity}
\Longleftrightarrow
\text{one-to-one correspondence}
\]

\section*{6. Constructing the Natural Numbers}

Frege defines zero using an impossible concept.

\[
Zx \equiv x \neq x
\]

Since nothing is non-self-identical, \(Z\) applies to no objects.

\[
0 = \#Z
\]

Then one is the number of the concept ``identical with zero'':

\[
1 = \#(x=0)
\]

Two is the number of the concept ``identical with zero or identical with one'':

\[
2 = \#(x=0 \lor x=1)
\]

And so on.

\subsection*{Teaching Point}

Frege tries to generate the natural numbers from logical notions: object, concept, identity,
and equinumerosity.

\subsection*{Whiteboard}

\[
0 = \#(x \neq x)
\]

\[
1 = \#(x=0)
\]

\[
2 = \#(x=0 \lor x=1)
\]

\[
3 = \#(x=0 \lor x=1 \lor x=2)
\]

\section*{7. Successor and Natural Number}

Frege defines the successor relation without presupposing arithmetic.

Roughly:

\[
n \text{ is the successor of } m
\]

if there is a concept \(F\) such that \(F\) applies to exactly \(n\) objects, and removing one
object from \(F\) leaves exactly \(m\) objects.

\[
n = m+1
\]

if \(n\) is the number of a concept obtained by adding one object to a concept with number
\(m\).

Frege then defines natural number without circularly appealing to ``finitely many.''

\[
n \text{ is natural iff } n \text{ belongs to every concept containing }0
\text{ and closed under successor.}
\]

\subsection*{Teaching Point}

Frege's definition of natural number anticipates the structure of mathematical induction.

\subsection*{Whiteboard}

\[
0 \in F
\]

\[
\forall n(n \in F \rightarrow S(n) \in F)
\]

\[
\therefore
\text{every natural number belongs to }F
\]

\section*{8. Frege's Theorem}

Frege derives the basic laws of arithmetic from Hume's Principle and second-order logic.

This result is now called Frege's Theorem.

\[
\boxed{
\text{Hume's Principle} + \text{second-order logic}
\Rightarrow
\text{arithmetic}
}
\]

The theorem is a major mathematical and philosophical achievement.

\subsection*{Teaching Point}

Even if Frege's full system fails, Frege's Theorem shows that arithmetic has a surprisingly
deep connection with logic and abstraction.

\subsection*{Whiteboard}

\[
\text{Hume's Principle}
\]

\[
+
\]

\[
\text{Second-order logic}
\]

\[
\Downarrow
\]

\[
\text{Arithmetic}
\]

\section*{9. The Caesar Problem}

Frege was not satisfied with Hume's Principle alone.

Hume's Principle settles identities of this form:

\[
\#F = \#G
\]

But it does not settle identities of this form:

\[
\#F = t
\]

where \(t\) is an arbitrary singular term.

For example:

\[
\#F = \text{Julius Caesar?}
\]

This is the Caesar problem.

The question is not whether anyone will actually confuse the number \(2\) with Julius
Caesar. The question is whether Frege has given full identity conditions for numbers as
objects.

\subsection*{Teaching Point}

If numbers are objects, we need to know when a number is identical with, or distinct from,
any object whatsoever.

\subsection*{Whiteboard}

\[
\#F = \#G
\quad
\text{settled}
\]

\[
\#F = \text{Caesar}
\quad
\text{not settled}
\]

\[
\text{No entity without identity.}
\]

\section*{10. Extensions and Basic Law V}

Frege tries to solve the Caesar problem by defining numbers as extensions of concepts.

\[
\#F =
\text{the extension of the concept equinumerous with }F
\]

So the number two is, roughly, the extension containing all concepts that apply to exactly
two objects.

Frege's system depends on Basic Law V.

\[
\boxed{
\operatorname{Ext}(F)=\operatorname{Ext}(G)
\quad \text{iff}
\quad
\forall x(Fx \leftrightarrow Gx)
}
\]

In words: two concepts have the same extension if and only if they apply to exactly the
same objects.

\subsection*{Teaching Point}

Frege's logicism depends on a theory of extensions. That theory is where disaster strikes.

\subsection*{Whiteboard}

\[
\operatorname{Ext}(F)=\operatorname{Ext}(G)
\Longleftrightarrow
\forall x(Fx \leftrightarrow Gx)
\]

\section*{11. Russell's Paradox}

Russell discovered that Basic Law V is inconsistent.

A simplified version of the paradox uses the class:

\[
R = \{x : x \notin x\}
\]

Now ask:

\[
R \in R?
\]

If \(R \in R\), then by definition \(R \notin R\).

If \(R \notin R\), then by definition \(R \in R\).

So:

\[
R \in R
\Longleftrightarrow
R \notin R
\]

Contradiction.

\subsection*{Teaching Point}

Russell's paradox destroyed Frege's original logicist system, because it showed that
Frege's theory of extensions was inconsistent.

\subsection*{Whiteboard}

\[
R = \{x : x \notin x\}
\]

\[
R \in R?
\]

\[
\text{Yes} \Rightarrow \text{No}
\]

\[
\text{No} \Rightarrow \text{Yes}
\]

\section*{12. Russell's Response: The Vicious Circle Principle}

Russell diagnoses the paradox as the result of illegitimate circularity.

\[
\boxed{
\text{No totality may contain members defined in terms of that very totality.}
}
\]

This is the vicious circle principle.

It rules out impredicative definitions: definitions that define something by referring to a
totality that already contains it.

\subsection*{Teaching Point}

Russell's response to paradox is to restrict what can meaningfully be defined.

\subsection*{Whiteboard}

\[
\text{Impredicative definition}
\]

\[
=
\]

\[
\text{definition using a totality containing the thing defined}
\]

\[
\text{Russell: ban it.}
\]

\section*{13. Type Theory}

Russell introduces type theory to block self-membership and vicious circularity.

\[
\begin{array}{c|c}
\textbf{Type 0} & \text{individuals} \\
\textbf{Type 1} & \text{classes of individuals} \\
\textbf{Type 2} & \text{classes of classes of individuals} \\
\textbf{Type 3} & \text{classes of type 2 classes}
\end{array}
\]

A class cannot be a member of itself because it is always of a higher type than its members.

\[
x \in C
\]

is meaningful only when \(C\) is one type higher than \(x\).

\subsection*{Teaching Point}

Type theory avoids Russell's paradox by making self-membership meaningless.

\subsection*{Whiteboard}

\[
\text{individuals}
\rightarrow
\text{classes}
\rightarrow
\text{classes of classes}
\rightarrow
\text{classes of classes of classes}
\]

\[
R \in R
\quad
\text{is not false; it is ill-typed.}
\]

\section*{14. Russell's Numbers}

Russell defines numbers as classes of equinumerous classes.

\[
2 = \text{the class of all pairs}
\]

\[
3 = \text{the class of all triples}
\]

So the number of a class \(C\) is the class of all classes equinumerous with \(C\).

This resembles Frege, but Russell's type restrictions create complications.

\subsection*{Teaching Point}

Russell preserves the general Fregean idea that numbers are tied to equinumerosity, but
he changes the logical framework to avoid paradox.

\subsection*{Whiteboard}

\[
0 = \text{class of all empty classes}
\]

\[
1 = \text{class of all singletons}
\]

\[
2 = \text{class of all pairs}
\]

\section*{15. The Axiom of Infinity}

Russell cannot prove that there are infinitely many natural numbers.

If there are only finitely many individuals, then there are only finitely many classes of
individuals, and the natural numbers run out.

So Russell and Whitehead introduce the axiom of infinity.

\[
\boxed{
\text{There are infinitely many individuals.}
}
\]

But this axiom does not look like a logical truth.

\subsection*{Teaching Point}

Russell avoids Frege's paradox, but he must add assumptions that weaken the logicist
claim.

\subsection*{Whiteboard}

\[
\text{Frege: proves infinity, but paradox.}
\]

\[
\text{Russell: avoids paradox, but assumes infinity.}
\]

\section*{16. The Axiom of Reducibility}

Ramified type theory also blocks many definitions needed for ordinary mathematics.

To recover mathematics, Russell and Whitehead introduce the axiom of reducibility.

Roughly:

\[
\boxed{
\text{For every class, there is an equivalent predicative class.}
}
\]

This allows them to proceed as if many impredicative definitions were acceptable.

But the axiom looks ad hoc. It is introduced because mathematics needs it, not because it
is obviously logical.

\subsection*{Teaching Point}

The axiom of reducibility reveals the cost of Russell's logicism: to save mathematics, he
must add principles that do not seem purely logical.

\subsection*{Whiteboard}

\[
\text{Avoid paradox}
\quad \rightarrow \quad
\text{type restrictions}
\]

\[
\text{Recover mathematics}
\quad \rightarrow \quad
\text{reducibility}
\]

\[
\text{Question: is this still logicism?}
\]

\section*{17. Russell's No-Class Theory}

Russell later denies that classes are genuine objects.

Classes are logical fictions.

Since numbers are constructed from classes, numbers also become logical fictions.

\[
\begin{array}{c|c}
\textbf{Frege} & \textbf{Russell} \\
\hline
\text{numbers are logical objects} & \text{numbers are logical fictions} \\
\text{realist ontology} & \text{anti-realist about classes} \\
\text{extensions as objects} & \text{classes eliminated by paraphrase}
\end{array}
\]

\subsection*{Teaching Point}

Frege and Russell are both logicists, but they differ sharply about what numbers are.

\subsection*{Whiteboard}

\[
\text{Frege: numbers exist as logical objects.}
\]

\[
\text{Russell: numbers are logical fictions.}
\]

\section*{18. Carnap and Logical Positivism}

Carnap shifts the debate from ontology to language.

He thinks the traditional question:

\[
\text{Do numbers really exist?}
\]

is either trivial or meaningless, depending on how it is asked.

Carnap distinguishes internal and external questions.

\[
\begin{array}{c|c}
\textbf{Internal questions} & \textbf{External questions} \\
\hline
\text{asked within a framework} & \text{asked about the framework as a whole} \\
\text{answered by rules} & \text{pragmatic or pseudo-question} \\
\text{Is there a prime greater than 100?} & \text{Do numbers really exist?}
\end{array}
\]

\subsection*{Teaching Point}

For Carnap, mathematics is analytic within a linguistic framework. The metaphysical
question of whether numbers really exist is not a genuine factual question.

\subsection*{Whiteboard}

\[
\text{Internal:}
\quad
\exists n(n \text{ is prime and } n>100)
\]

\[
\text{External:}
\quad
\text{Do numbers really exist?}
\]

\[
\text{Carnap: choose frameworks pragmatically.}
\]

\section*{19. Carnap's Principle of Tolerance}

Carnap defends a principle of tolerance.

We are free to adopt whatever linguistic framework is useful, provided the rules are clear.

The question is not:

\[
\text{Is this framework metaphysically correct?}
\]

The question is:

\[
\text{Is this framework useful for science and inquiry?}
\]

This lets Carnap accept impredicative systems if they are useful, without defending them
by metaphysics.

\subsection*{Teaching Point}

Carnap replaces metaphysical justification with pragmatic choice of language.

\subsection*{Whiteboard}

\[
\text{Frameworks are tools.}
\]

\[
\text{Choose by clarity, efficiency, and usefulness.}
\]

\[
\text{No metaphysical tribunal.}
\]

\section*{20. Ayer and Mathematical Truth}

Ayer states the positivist view clearly.

Empiricism says that factual truths are known through experience.

But mathematics seems necessary and certain.

So the positivist concludes:

\[
\boxed{
\text{Mathematics is necessary because it has no factual content.}
}
\]

Mathematical truths are analytic. They are true by virtue of linguistic rules.

\[
\begin{array}{c|c}
\textbf{Science} & \textbf{Mathematics} \\
\hline
\text{factual} & \text{analytic} \\
\text{empirically testable} & \text{true by rules} \\
\text{fallible} & \text{necessary}
\end{array}
\]

\subsection*{Teaching Point}

Logical positivism preserves empiricism by saying that mathematics is not factual.

\subsection*{Whiteboard}

\[
\text{If factual, then empirical.}
\]

\[
\text{Mathematics is not empirical.}
\]

\[
\therefore
\quad
\text{mathematics is not factual.}
\]

\section*{21. Geometry as a Framework}

Logical positivists also treat geometry as a framework.

Euclidean geometry is not refuted by measuring a triangle and finding that its angles do
not sum to \(180^\circ\).

Instead, we say:

\[
\text{the measurement was wrong}
\]

or

\[
\text{the physical triangle is not Euclidean.}
\]

Pure geometry is analytic within its framework.

The empirical question is which geometry is useful for physics.

\subsection*{Teaching Point}

The logical positivists separate pure mathematics from applied mathematics.

\subsection*{Whiteboard}

\[
\text{Pure geometry: analytic framework}
\]

\[
\text{Applied geometry: empirical choice}
\]

\[
\text{Euclidean? Non-Euclidean?}
\]

\section*{22. Problems for Logical Positivism}

Logical positivism eventually runs into several major problems.

\[
\begin{array}{c|p{0.62\textwidth}}
\textbf{Problem} & \textbf{Question} \\
\hline
\text{Verification} &
What exactly makes a factual statement verifiable? \\

\text{Self-reference} &
Is the verification principle itself analytic or empirically verifiable? \\

\text{Analytic/synthetic distinction} &
Can we sharply distinguish truths by meaning from truths by fact? \\

\text{Mathematical practice} &
Is mathematics really divided into self-contained frameworks?
\end{array}
\]

Quine attacks the analytic/synthetic distinction and develops a more holistic empiricism.

\subsection*{Teaching Point}

The positivist account of mathematics depends on a sharp distinction between analytic
framework rules and empirical factual claims. That distinction comes under pressure.

\subsection*{Whiteboard}

\[
\text{Analytic}
\quad \big| \quad
\text{Synthetic}
\]

\[
\text{Quine: no sharp boundary.}
\]

\section*{23. G\"odel and Self-Contained Frameworks}

Carnap's picture suggests that mathematical truths are knowable from the rules of their
framework.

But G\"odel's incompleteness theorem complicates this.

If \(D\) is a sufficiently strong effective deductive system, then there are sentences in the
language of \(D\) that are not decided by the rules of \(D\).

\[
D \nvdash P
\qquad
D \nvdash \neg P
\]

In practice, mathematicians often prove results about one structure by embedding it in a
richer structure.

\subsection*{Teaching Point}

Mathematical knowledge is not neatly confined inside isolated linguistic frameworks.

\subsection*{Whiteboard}

\[
\text{Arithmetic statement}
\]

\[
\Downarrow
\]

\[
\text{proved using richer mathematics}
\]

\[
\text{Frameworks are not self-contained.}
\]

\section*{24. Fermat's Last Theorem}

Fermat's Last Theorem is a useful example.

It is a statement about natural numbers:

\[
a^n+b^n \neq c^n
\]

for positive natural numbers \(a,b,c\) and \(n>2\).

The statement is elementary enough to understand with basic arithmetic.

But its proof uses highly sophisticated mathematics far beyond elementary arithmetic.

\subsection*{Teaching Point}

Understanding the statement of a theorem is not the same as having the resources to know
that it is true.

\subsection*{Whiteboard}

\[
\text{Simple statement}
\]

\[
a^n+b^n \neq c^n
\]

\[
\text{Deep proof}
\]

\[
\text{elliptic curves, modular forms, advanced structures}
\]

\section*{25. Neo-Logicism}

Contemporary neo-logicists, such as Crispin Wright and Bob Hale, revive Frege's project.

They avoid Basic Law V and work directly with Hume's Principle.

\[
\boxed{
\#F = \#G
\quad \text{iff}
\quad
F \sim G
}
\]

The neo-logicist claim is not usually that Hume's Principle is a logical truth in the strict
old sense.

Rather, Hume's Principle is treated as analytic, or meaning-constitutive, of the concept of
cardinal number.

\[
\text{Hume's Principle}
+
\text{second-order logic}
\Rightarrow
\text{arithmetic}
\]

\subsection*{Teaching Point}

Neo-logicism preserves Frege's central insight while avoiding the inconsistent theory of
extensions.

\subsection*{Whiteboard}

\[
\text{Frege: Basic Law V}
\quad \Rightarrow \quad
\text{paradox}
\]

\[
\text{Neo-logicism: Hume's Principle}
\quad \Rightarrow \quad
\text{arithmetic}
\]

\section*{26. Problems for Neo-Logicism}

Neo-logicism faces several major problems.

\[
\begin{array}{c|p{0.62\textwidth}}
\textbf{Problem} & \textbf{Question} \\
\hline
\text{Second-order logic} &
Is second-order logic really logic, or is it set theory in disguise? \\

\text{Caesar problem} &
Does Hume's Principle fully identify numbers as objects? \\

\text{Bad company} &
How do we distinguish good abstraction principles from bad ones? \\

\text{Scope} &
Can the method extend beyond arithmetic?
\end{array}
\]

The bad company objection is especially important.

Basic Law V and Hume's Principle are both abstraction principles. One is inconsistent; the
other is fruitful.

\[
\text{What separates good abstraction from bad abstraction?}
\]

\subsection*{Teaching Point}

Neo-logicism depends on showing that Hume's Principle is legitimate in a way that bad
abstraction principles are not.

\subsection*{Whiteboard}

\[
\text{Good abstraction: Hume's Principle}
\]

\[
\text{Bad abstraction: Basic Law V}
\]

\[
\text{Need a non-ad hoc distinction.}
\]

\section*{Conclusion: What Would Logicism Have to Show?}

Logicism must show that mathematics can be derived from logic and definitions.

But the chapter repeatedly presses the same question:

\[
\boxed{
\text{What counts as logic?}
}
\]

If logic includes concepts, extensions, classes, second-order quantification, infinity,
choice, reducibility, and abstraction principles, then the reduction may look suspicious.

Perhaps mathematics has not been reduced to logic.

Perhaps mathematics has been moved into logic.

\[
\begin{array}{c|c|c}
\textbf{Frege} & \textbf{Russell} & \textbf{Carnap} \\
\hline
\text{numbers as logical objects} &
\text{numbers as logical fictions} &
\text{numbers as framework-internal entities} \\
\text{realist logicism} &
\text{type-theoretic logicism} &
\text{linguistic logicism} \\
\text{paradox problem} &
\text{extra axioms problem} &
\text{framework problem}
\end{array}
\]

\section*{Final Framing}

To teach this chapter well, present logicism as a powerful but unstable answer to the
post-Kantian problem.

It gives a beautiful explanation of the a priori character of mathematics:

\[
\text{Mathematics is knowable because it follows from logic.}
\]

But it also faces a deep challenge:

\[
\text{The more logic must contain, the less clear it is that mathematics has been reduced.}
\]

\[
\boxed{
\text{Logicism lives or dies by the boundary between logic and mathematics.}
}
\]

\section*{Discussion Questions}

\begin{enumerate}[leftmargin=2em]
    \item What does logicism claim about mathematics?
    \item How does Frege's notion of analyticity differ from Kant's?
    \item Why is equinumerosity important for defining number?
    \item What does Hume's Principle say?
    \item Does Hume's Principle fully identify numbers as objects?
    \item What is the Caesar problem?
    \item Why does Russell's paradox undermine Frege's system?
    \item Does type theory solve the paradox at too high a cost?
    \item Is the axiom of infinity a logical truth?
    \item Why does the axiom of reducibility look ad hoc?
    \item Are numbers objects, fictions, or framework-internal entities?
    \item Is Carnap right that external ontological questions are pseudo-questions?
    \item Does G\"odel's incompleteness theorem undermine the idea of self-contained mathematical frameworks?
    \item Is second-order logic really logic?
    \item Can neo-logicism succeed where Frege's original logicism failed?
\end{enumerate}

\section*{Key Terms}

\begin{description}[leftmargin=2em]
    \item[Logicism] The view that mathematics, or a significant part of mathematics, can be reduced to logic.

    \item[Analytic truth] For Frege, a truth provable from logic and definitions alone.

    \item[A priori] Knowable independently of sensory experience.

    \item[Concept] For Frege, something under which objects fall; roughly, a property.

    \item[Equinumerosity] The relation between two concepts or collections when their objects can be placed in one-to-one correspondence.

    \item[Hume's Principle] The principle that the number of \(F\) is identical with the number of \(G\) if and only if \(F\) and \(G\) are equinumerous.

    \item[Frege's Theorem] The result that arithmetic can be derived from Hume's Principle in second-order logic.

    \item[Caesar problem] The problem that Hume's Principle determines identities between numbers but not identities between numbers and arbitrary objects.

    \item[Extension] The class of all objects falling under a concept.

    \item[Basic Law V] Frege's principle that two concepts have the same extension if and only if they apply to the same objects.

    \item[Russell's paradox] The contradiction generated by considering the class of all classes that are not members of themselves.

    \item[Impredicative definition] A definition that defines an object by referring to a totality containing that object.

    \item[Vicious circle principle] Russell's principle that no object may be defined by reference to a totality containing that object.

    \item[Type theory] Russell's hierarchy of individuals, classes of individuals, classes of classes, and so on, designed to avoid paradox.

    \item[Axiom of infinity] Russell and Whitehead's postulate that there are infinitely many individuals.

    \item[Axiom of reducibility] Russell and Whitehead's postulate that every class is equivalent to a predicative class.

    \item[Logical fiction] Russell's term for entities, such as classes, that are useful in discourse but not part of the ultimate furniture of reality.

    \item[Linguistic framework] Carnap's term for a system of language and rules within which questions can be meaningfully asked and answered.

    \item[Internal question] For Carnap, a question asked within a linguistic framework.

    \item[External question] For Carnap, a question about the reality or legitimacy of a framework as a whole.

    \item[Principle of tolerance] Carnap's view that we may choose whichever linguistic frameworks are useful, provided their rules are clear.

    \item[Neo-logicism] A contemporary revival of Fregean logicism, especially through Hume's Principle and abstraction principles.

    \item[Bad company objection] The objection that neo-logicists need a principled way to distinguish acceptable abstraction principles from unacceptable ones.
\end{description}

\end{document}